%% v3 copy of table_molecules.tex with the two strings Sec. 9.2 retracts removed.
%% The repository file is NOT edited by this session; this copy must replace it.
\begin{table*}[tp]\centering\small
\caption{\textbf{The nineteen-molecule fermionic-magic suite}, computed in a cc-pVDZ
active space; the active-space CASCI energies are converged to $1.1\times10^{-13}$~Ha, which is a
convergence check and not a validation against an independent method (Sec.~\ref{sec:scope:mol}). The fermionic
AntiFlatness $\mathcal F_1$ ($k{=}1$ Majorana antiflatness; $\mathcal F_1{=}0$ iff Gaussian/mean-field, maximum
$\mathcal F_1^{\max}{=}2n_{\rm o}$) is listed near equilibrium and at dissociation, in a valence
bond-breaking active space CAS$(n_e,n_{\rm o})$. Fermionic magic switches on for covalent
bond-breaking, is already large at equilibrium for inherently multireference species
($\mathrm{C_2}$, $\mathrm{O_2}$), and stays near-Gaussian for ionic and weakly-correlated
bonds---so $\mathcal F_1$ marks the multireference, strongly-correlated regime rather than bond length
\emph{per se}. The last two columns validate the \emph{method}: at the dissociation geometry the
reconstruction of the (orbital-projected) spectral function $A(\omega)$ from a subspace drawn in
$2.16\times10^{5}$ shots matches the exact Lehmann result to relative-$L_1$ error, on a subspace that
is a fraction $|\mathcal{S}|/D$ of the full $(N{+}1)$ sector. In twelve of the nineteen the subspace contains the probe's Krylov space and the agreement is exact by construction; see
Table~\ref{tab:molscope}. The listed $\mathcal F_1=4\,\mathrm{tr}[\gamma(1-\gamma)]$
are directly computed spin-orbital values (for the five open-shell species they are \emph{not} $2N_{\rm u}$;
see Sec.~\ref{sec:nogo}). Per-molecule equilibrium and dissociation geometries, active-space orbital
indices, FCI energies and residuals, and the absolute $|\mathcal S|$, sector dimension $D$ and minimal bond
dimension $\chi$ are listed in Sec.~\ref{app:repro}. $^{\dagger}$: near-degenerate ground state.}
\label{tab:molecules}
\begin{tabular}{llcccccc}\hline\hline
Molecule & Bonding & CAS$(n_e,n_{\rm o})$ & $\mathcal F_1^{\rm eq}$ & $\mathcal F_1^{\rm diss}$ & $\mathcal F_1^{\rm diss}/\mathcal F_1^{\max}$ & rel-$L_1$ & $|\mathcal{S}|/D$ \\ \hline
\multicolumn{8}{l}{\emph{Covalent bond-breaking}}\\
F$_2$ & covalent single & (1+1,2) & 0.068 & 3.985 & 1.00 & $<\!10^{-13}$ & 0.50 \\
N$_2$ & covalent triple & (3+3,6) & 0.847 & 11.377 & 0.95 & $2\times10^{-5}$ & 0.21 \\
H$_2$O & multicentre & (2+2,4) & 0.007 & 6.936 & 0.87 & $<\!10^{-13}$ & 0.50 \\
HF & polar single & (1+1,2) & 0.001 & 3.460 & 0.87 & $<\!10^{-13}$ & 1.00 \\
NH$_3$ & multicentre & (3+3,6) & 0.080 & 10.289 & 0.86 & $1\times10^{-5}$ & 0.47 \\
CO & polar triple & (3+3,6) & 0.688 & 8.192 & 0.68 & $1\times10^{-4}$ & 0.41 \\
H$_6$ & H-chain & (3+3,6) & 0.521 & 7.735 & 0.64 & $2\times10^{-5}$ & 0.50 \\
H$_4$ & H-chain & (2+2,4) & 0.252 & 5.113 & 0.64 & $<\!10^{-13}$ & 0.50 \\
H$_2$ & covalent single & (1+1,2) & 0.038 & 2.511 & 0.63 & $<\!10^{-13}$ & 0.50 \\
BeH$_2$ & multicentre & (2+2,4) & 0.039 & 3.754 & 0.47 & $<\!10^{-13}$ & 0.25 \\
\multicolumn{8}{l}{\emph{Multireference / open-shell}}\\
C$_2$$^{\dagger}$ & covalent double & (4+4,6) & 4.743 & 8.001 & 0.67 & $<\!10^{-13}$ & 0.10 \\
NO & polar (radical) & (4+3,6) & 0.790 & 7.449 & 0.62 & $2\times10^{-6}$ & 0.25 \\
O$_2$ & covalent double (triplet) & (5+3,6) & 1.673 & 6.000 & 0.50 & $1\times10^{-5}$ & 0.25 \\
CN & polar (radical) & (4+3,6) & 1.307 & 5.938 & 0.49 & $1\times10^{-5}$ & 0.24 \\
OH$^{\dagger}$ & polar (radical) & (3+2,4) & 0.013 & 3.703 & 0.46 & $<\!10^{-13}$ & 0.33 \\
\multicolumn{8}{l}{\emph{Ionic / weakly correlated}}\\
BH & metal hydride & (2+2,4) & 0.518 & 0.573 & 0.07 & $<\!10^{-13}$ & 0.25 \\
BeH & metal hydride (radical) & (2+1,3) & 0.035 & 0.386 & 0.06 & $<\!10^{-13}$ & 0.67 \\
LiH & metal hydride & (1+1,2) & 0.002 & 0.075 & 0.02 & $<\!10^{-13}$ & 1.00 \\
LiF & ionic & (1+1,2) & 0.000 & 0.001 & 0.00 & $<\!10^{-13}$ & 1.00 \\
\hline\hline\end{tabular}
\end{table*}
